\begin{frame}
  \frametitle{Simplification of the \gls{bte}}
  Because the \glspl{arp} are stated to be \textit{non-multiplying}, \textit{time-independent}, and \textit{one-speed}, \cref{eq:general-bte} is immediately simplified to:
  \begin{multline}\label{eq:first-simplification-bte}
    \ohat \cdot \grad \psisimple + \Sigma_t \left(\rvec\right)\psisimple =\\=\int_{4\pi}\sum_{m=0}^{\infty}\frac{2m+1}{4\pi}P_m \left(\ohat' \cdot \ohat \right)\Sigma_{s,m}\left(\rvec\right)\psi \left(\rvec,\ohat'\right) \dd \ohat' + \frac{Q\left(\rvec\right)}{4\pi}
  \end{multline}
  Note that we have also substituted in \cref{eq:scattering-poly-expansion}.
\end{frame}

\begin{frame}
  \frametitle{Simplification of the \gls{bte}}
  The use of constant-value properties has a couple of immediate consequences, notably in the inscatter term.
  For a constant-value isotropic scattering cross section, \cref{eq:legendre-moments} is readily evaluated:
  \begin{equation}\label{eq:inscatter-xs-evaluated}
    \begin{split}
      \Sigma_{s,m}\left(\rvec\right) & = 2\pi \int_{-1}^{1}\Sigma_s \left(\rvec,\mu_0\right)P_m \left(\mu_0\right)\dd \mu_0 \\
                                     & =2\pi \int_{-1}^{1}\frac{\Sigma_s}{4\pi} P_m \left(\mu_0\right)\dd \mu_0             \\
                                     & =\frac{\Sigma_s}{2}\int_{-1}^{1} P_m \left(\mu_0\right)\dd \mu_0                     \\
                                     & =\Sigma_s \delta_{0m}
    \end{split}
  \end{equation}
  where, in the last line, the orthogonality of Legendre polynomials is applied.

\end{frame}

\begin{frame}
  \frametitle{Simplification of the \gls{bte}}
  Substitution of \cref{eq:inscatter-xs-evaluated} into \cref{eq:first-simplification-bte} gives:
  \begin{equation}
    \ohat \cdot \grad \psisimple + \Sigma_t \left(\rvec\right)\psisimple =\frac{\Sigma_s}{4\pi}\int_{4\pi}\psi \left(\rvec,\ohat'\right) \dd \ohat' + \frac{Q\left(\rvec\right)}{4\pi}
  \end{equation}
  Or:
  \begin{equation} \label{eq:fully-simplified-bte}
    \ohat \cdot \grad \psisimple + \Sigma_t \left(\rvec\right)\psisimple =\frac{\Sigma_s}{4\pi}\phi\left(\rvec\right) + \frac{Q\left(\rvec\right)}{4\pi}
  \end{equation}
  where:
  \begin{equation}
    \label{eq:scalar-flux}
    \phi\left(\rvec\right)\equiv \int_{4\pi}\psi \left(\rvec,\ohat'\right) \dd \ohat'
  \end{equation}
  \Cref{eq:scalar-flux} is know as the \textit{angle-integrated} (or \textit{scalar}) \textit{flux}.
\end{frame}

